\lecture{24}{Двойственный подъём и декомпозиция}

\sourcevideo
    {https://www.youtube.com/playlist?list=PLOzRYVm0a65fhKDS8cYIC7YCuVGJ4FOJx}
    {Week 7: Lecture 24: Dual Ascent and Dual Decomposition}

Рассматривается выпуклая задача с линейным равенством
\[
    \min_{x\in\mathbb R^n} f(x)
    \qquad
    \text{при условии}
    \qquad
    Ax=b.
\]
Двойственный подъём вычисляет градиент двойственной функции через
вспомогательную минимизацию по прямой переменной. Если функция и
переменные разделяются на блоки, эта минимизация распадается на
независимые локальные задачи.

\section{Сопряжённая функция и субградиенты}

Для собственной функции
\[
    f\colon\mathbb R^n\to\mathbb R\cup\{+\infty\}
\]
сопряжённая функция Фенхеля определяется формулой
\[
    f^*(y)
    =
    \sup_{x\in\mathbb R^n}
    \bigl\{\langle y,x\rangle-f(x)\bigr\}.
\]
Она всегда выпукла как поточечный супремум аффинных функций.
Эпиграф функции равен
\[
    \operatorname{epi}f
    =
    \bigl\{(x,t)\in\mathbb R^{n+1}:f(x)\le t\bigr\}.
\]
Функция называется замкнутой, если её эпиграф замкнут; для собственной
выпуклой функции это эквивалентно нижней полунепрерывности.

Повторное сопряжение задаётся формулой
\[
    f^{**}(x)
    =
    \sup_y\bigl\{\langle x,y\rangle-f^*(y)\bigr\}.
\]
По теореме Фенхеля--Моро для собственной замкнутой выпуклой функции
\[
    f^{**}=f.
\]
В общем случае \(f^{**}\) является замкнутой выпуклой оболочкой
исходной функции.

Из определения сопряжения следует неравенство Фенхеля--Юнга
\[
    f(x)+f^*(y)\ge \langle x,y\rangle.
\]
Для собственной замкнутой выпуклой функции эквивалентны условия
\[
    y\in\partial f(x),
    \qquad
    x\in\partial f^*(y),
    \qquad
    f(x)+f^*(y)=\langle x,y\rangle.
\]
При дифференцируемости обеих функций это даёт
\[
    y=\nabla f(x)
    \quad\Longleftrightarrow\quad
    x=\nabla f^*(y).
\]

\section{Двойственная функция}

Пусть
\[
    A\in\mathbb R^{m\times n},
    \qquad
    b\in\mathbb R^m.
\]
Функция Лагранжа и двойственная функция имеют вид
\[
    \mathcal L(x,\nu)
    =
    f(x)+\nu^{\mathsf T}(Ax-b),
    \qquad
    g(\nu)=\inf_x\mathcal L(x,\nu).
\]
Используя определение сопряжённой функции, получаем
\[
\begin{aligned}
    g(\nu)
    &=
    -b^{\mathsf T}\nu
    +
    \inf_x
    \bigl\{f(x)+\langle A^{\mathsf T}\nu,x\rangle\bigr\}
    \\
    &=
    -f^*(-A^{\mathsf T}\nu)-b^{\mathsf T}\nu.
\end{aligned}
\]
Двойственная задача состоит в максимизации \(g\) по
\(\nu\in\mathbb R^m\).

Предположим, что для каждого \(\nu\) минимум достигается в единственной
точке
\[
    x(\nu)
    =
    \operatorname*{argmin}_x
    \bigl\{f(x)+\nu^{\mathsf T}(Ax-b)\bigr\}.
\]
По теореме Данскина
\[
    \nabla g(\nu)=Ax(\nu)-b.
\]
Таким образом, двойственный градиент совпадает с невязкой прямого
ограничения. Явную формулу для \(f^*\) вычислять не требуется:
достаточно решить задачу минимизации по \(x\).

При нескольких минимизаторах функция \(g\) может быть
недифференцируемой. Для любого
\[
    x\in\operatorname*{argmin}_z\mathcal L(z,\nu)
\]
вектор \(Ax-b\) является суперградиентом вогнутой функции \(g\).
Сильная выпуклость \(f\) гарантирует единственность \(x(\nu)\) и
дифференцируемость двойственной функции.

\section{Двойственный подъём и скорость сходимости}

Одна итерация двойственного подъёма состоит из прямой минимизации
\[
    x^{k+1}
    =
    \operatorname*{argmin}_x
    \bigl\{f(x)+(\nu^k)^{\mathsf T}(Ax-b)\bigr\}
\]
и обновления
\[
    \nu^{k+1}
    =
    \nu^k+\alpha_k\bigl(Ax^{k+1}-b\bigr),
    \qquad
    \alpha_k>0.
\]
Поскольку максимизируется вогнутая функция, это именно подъём; запись
как спуска относится к минимизации \(-g\).

Если \(f\) является \(\mu\)-сильно выпуклой, то \(f^*\) имеет
\(1/\mu\)-липшицев градиент. Поэтому
\[
    \lVert\nabla g(\nu_1)-\nabla g(\nu_2)\rVert
    \le
    \frac{\lVert A\rVert^2}{\mu}
    \lVert\nu_1-\nu_2\rVert.
\]
Обозначим
\[
    L_d=\frac{\lVert A\rVert^2}{\mu}.
\]
Для шага \(\alpha=1/L_d\) и существующего двойственного решения
\(\nu^*\) стандартная оценка гладкого градиентного подъёма даёт
\[
    g^*-g(\nu^k)
    \le
    \frac{L_d\lVert\nu^0-\nu^*\rVert^2}{2k}.
\]
Следовательно, достижение двойственной точности \(\varepsilon\)
требует порядка \(O(\varepsilon^{-1})\) итераций.

Если \(f\) дополнительно \(L\)-гладка, а \(A\) имеет полный строчный
ранг, то \(-g\) сильно выпукла с константой
\[
    \mu_d
    =
    \frac{\sigma_{\min}(A)^2}{L}.
\]
При шаге
\[
    \alpha=\frac{2}{L_d+\mu_d}
\]
получаем линейное сжатие
\[
    \lVert\nu^{k+1}-\nu^*\rVert
    \le
    \frac{L_d-\mu_d}{L_d+\mu_d}
    \lVert\nu^k-\nu^*\rVert.
\]
Двойственное число обусловленности равно
\[
    \kappa_d
    =
    \frac{L_d}{\mu_d}
    =
    \frac{L\lVert A\rVert^2}
         {\mu\sigma_{\min}(A)^2},
\]
поэтому скорость зависит и от геометрии функции, и от обусловленности
линейного ограничения.

На практике прямая подзадача часто решается приближённо. Если вместо
\(x(\nu^k)\) используется \(\widetilde x^{k+1}\), то обновление
строится по вектору
\[
    A\widetilde x^{k+1}-b.
\]
Для сохранения сходимости погрешности внутренних решений должны
убывать или быть суммируемыми; точное условие зависит от выбранной
внешней схемы.

\section{Блочно-разделимая задача}

Рассмотрим задачу
\[
\begin{aligned}
    \min_{x_1,\ldots,x_B}\quad
    &\sum_{i=1}^{B}f_i(x_i),
    \\
    \text{при условии}\quad
    &\sum_{i=1}^{B}A_ix_i=b,
\end{aligned}
\]
где
\[
    x_i\in\mathbb R^{n_i},
    \qquad
    A_i\in\mathbb R^{m\times n_i}.
\]
Её функция Лагранжа разделяется:
\[
    \mathcal L(x_1,\ldots,x_B,\nu)
    =
    \sum_{i=1}^{B}
    \bigl(f_i(x_i)+\nu^{\mathsf T}A_ix_i\bigr)
    -\nu^{\mathsf T}b.
\]
Поэтому при фиксированном \(\nu\) минимизация по блокам независима,
а двойственная функция равна
\[
    g(\nu)
    =
    -\nu^{\mathsf T}b
    -
    \sum_{i=1}^{B}f_i^*(-A_i^{\mathsf T}\nu).
\]

На итерации \(k\) каждый агент параллельно вычисляет
\[
    x_i^{k+1}
    =
    \operatorname*{argmin}_{x_i}
    \bigl\{f_i(x_i)+(\nu^k)^{\mathsf T}A_ix_i\bigr\}.
\]
После этого координатор формирует невязку
\[
    r^{k+1}
    =
    \sum_{i=1}^{B}A_ix_i^{k+1}-b
\]
и обновляет общий множитель
\[
    \nu^{k+1}
    =
    \nu^k+\alpha_k r^{k+1}.
\]
Именно распад прямой минимизации после перехода к двойственной задаче
называется двойственной декомпозицией.

Разделимость нарушается, если целевая функция содержит член
\(\psi(x_i,x_j)\), связывающий разные блоки. Тогда связанные
переменные объединяют в один блок, вводят локальные копии с условиями
согласования или применяют другой метод декомпозиции. Например,
\[
    \frac12x_1^2+\frac12x_2^2+x_2x_3
    =
    f_1(x_1)+f_2(x_2,x_3),
\]
поэтому \(x_2\) и \(x_3\) образуют один блок. Сам этот пример
иллюстрирует только структуру разделения: без дополнительных условий
локальная функция не обязана быть выпуклой.

\section{Архитектура и коммуникации}

Алгоритм имеет схему «рассылка--сбор». Координатор передаёт всем
агентам \(\nu^k\), агенты параллельно решают локальные задачи и
возвращают векторы \(A_ix_i^{k+1}\), после чего координатор вычисляет
общую невязку и новый множитель. При ограничении размерности \(m\)
объём передачи в каждом направлении имеет порядок \(O(Bm)\) скалярных
чисел за синхронную итерацию. Размер локальной переменной \(n_i\) не
обязан совпадать с размером сообщения.

Передача только \(A_ix_i\), а не полной функции \(f_i\) или вектора
\(x_i\), уменьшает объём раскрываемой информации, но не обеспечивает
конфиденциальность. При известной матрице \(A_i\) и малом ядре
локальная переменная может быть частично восстановлена; повторные
сообщения также несут информацию о функции. Формальные гарантии
требуют дополнительных механизмов защиты.

Несмотря на параллельные вычисления, схема остаётся централизованной:
координатор агрегирует сообщения, обновляет множитель и является
единой точкой отказа. Синхронная версия ждёт самого медленного агента;
асинхронные варианты допускают устаревшие сообщения, но требуют
отдельного анализа. Двойственную переменную можно интерпретировать как
цену общего ресурса: положительный дисбаланс увеличивает цену, а
недоиспользование уменьшает её.

\section{Связь с дополненной функцией Лагранжа}

Обычный двойственный подъём использует
\[
    \mathcal L(x,\nu)
    =
    f(x)+\nu^{\mathsf T}(Ax-b).
\]
Метод множителей заменяет её дополненной функцией
\[
    \mathcal L_c(x,\nu)
    =
    f(x)+\nu^{\mathsf T}(Ax-b)
    +
    \frac c2\lVert Ax-b\rVert^2.
\]
Квадратичный штраф улучшает свойства двойственного шага, но для
блочного ограничения создаёт перекрёстные члены
\[
    x_i^{\mathsf T}A_i^{\mathsf T}A_jx_j,
    \qquad
    i\ne j,
\]
и тем самым обычно разрушает независимость локальных подзадач. ADMM
использует дополненную функцию Лагранжа, но минимизирует её по блокам
последовательно, сохраняя часть преимуществ декомпозиции.

Двойственная декомпозиция требует существования решений локальных
задач и подходящего выбора шага. При неединственности локальных
минимизаторов сходимость двойственных переменных не обязательно сразу
влечёт сходимость каждого прямого вектора. Условие консенсуса также
можно записать как линейное равенство, однако полностью
децентрализованные методы используют структуру графа и обходятся без
единого координатора.
